\documentclass[10pt]{article}
\usepackage{kotex}
\usepackage{amssymb}
\usepackage{amsmath}
\usepackage{graphicx}
\usepackage{subfigure}
%%%%%%%%%%%%%%%%%
\usepackage{titlesec}

\titleformat{\section} 
  {\centering\normalfont\large\sffamily\bfseries}
  {\thesection.}
  {1em}
  {}
\titleformat{\subsection} 
  {\normalfont\normalsize\sffamily\bf}
  {\thesubsection.}
  {1em}
  {\itshape}
%%%%%%%%%%%%%%%%%%
\renewcommand{\thesection}{\Roman{section}} 
\def\thesubsection{\arabic{subsection}}
 
%%%%%%%%%%%%%%%
\usepackage{setspace}
\setstretch{1.31}


\textwidth 17.6cm \textheight 24.2cm
\addtolength{\oddsidemargin}{-25mm}
\addtolength{\topmargin}{-21mm}
%\parskip 5mm
%\parindent 3mm
\setcounter{page}{0}
 
 
   

\begin{document}



\begin{center}
{ \LARGE \bf    논문 제목 }\\[12pt]
 학생. 지도교사 \\
 서울과학고등학교
\\[12pt]
{ \LARGE \bf    paper title }\\[12pt]
 author, teacher \\
 Seoul Science High School
  \\[12pt]
{\large \bf  국문 초록}
  \noindent\rule{\textwidth}{0.4pt}
%
  \begin{flushleft}  
  초록을 작성할 때는 연구 주제의 목적과 결과, 알아낸 사실들을 일목요연하게 정리하여 연구 전체를 하나의 문단으로 요약한다. 첫 줄은 들여쓰기하지 않으며 400단어 내외로 작성한다. ‘초록 내용’
주제어: 논문에서 다룬 5개 이내의 키워드. 제목에 사용된 단어가 포함해야 한다. ‘초록 내용’
  \noindent\rule{\textwidth}{0.4pt}
\end{flushleft}  
{\large \bf Abstract}
  \noindent\rule{\textwidth}{0.4pt}
    \begin{flushleft}  
When writing an abstract, summarize the entire research in one paragraph by arranging the purpose, results, and findings of the research topic at a glance. The first line is not indented and should be written within 400 words. ‘Abstract’
Keywords: List of 5 or less keywords covered in the paper. The words used in the title should be included.  ‘Abstract’
  \noindent\rule{\textwidth}{0.4pt}
\end{flushleft} 
  \end{center} 

 
%%%%%%%%%%%%%%%%%%%%%%%%%
\section{\large \bf 서론(Introduction) }
%%%%%%%%%%%%%%%%%%%%%%%%%
‘본문’스타일을을 사용한다. 강조할 단어에는 ‘굵게’스타일이나 ‘기울임’스타일을 사용한다. 
이 연구를 수행하게 된 이유와 목적, 연구의 동기와 필요성을 한두 문단으로 기술한다. 


%%%%%%%%%%%%%%%%%%%%%%
\section{\large \bf 이론적 배경(Theory \& Principle)}
%%%%%%%%%%%%%%%%%%%%%%
‘본문’스타일을 사용한다. 
이 연구가 시행되기 이전의 선행연구를 소개하여 연구 결과를 이해하기 전 알아야 할 사전 지식을 제공해야 한다. 
정리 2.1 ([인용논문,정리번호]) (선택사항) 본 연구와 관련된 정리를 서술한다.
정의 2.2 (선택사항) 본 연구와 관련된 정의를 서술한다.
 
 
 %%%%%%%%%%%%%%%%%%%%%%%
\section{\large \bf 연구 관련 제목(Chpater Name Related Research)}
%%%%%%%%%%%%%%%%%%%%%%%



%%%%%%%%%%%%%%%%%%%%%%%%%%%%%
\section{\large \bf  (선택사항) 결론 및 제언(Conclusion and Proposal)}
%%%%%%%%%%%%%%%%%%%%%%%%%%%%%
본문’스타일을 사용한다. 실험에서 연구의 가설을 설명하기 위해 시행했던 시험의 결과결과들을 보여주어야 하며, 각 실험들은 사진이나 그래프, 표 등으로 한눈에 알아볼 수 있도록 제시해주어야 한다.
%%%%%%%%%%%%%%%%%%%%%%%%%%%%%%%%%%%%%%%%%%%%%%%%%%%%% 
\begin{figure} [hbtp]
\centering
\includegraphics[scale=.4]{20240527_125844.jpg}
\caption{ Cross section  of a birdbath  type   optics system in the vertical direction }%
\label{fig2}
\end{figure}
%%%%%%%%%%%%%%%%%%%%%%%%%%%%%%%%%%%%%%%%%%%%%%%%%%%%% 




 
%%%%%%%%%%%%%%%%%%%%%%
\section*{\large \bf 보충 자료(Supplemental Information) }
%%%%%%%%%%%%%%%%%%%%%%
\appendix
\renewcommand{\theequation}{A.\arabic{equation}}
\setcounter{equation}{0}
{\subsection{\large \bf About quantum field}}
%%%%%%%%%%%%%%%%%%%%%%



%%%%%%%%%%%%%%%%%%%%%%%%%%%%%%%%%
\begingroup
\titleformat*{\section}  {\centering\normalfont\large\sffamily\bfseries}
\begin{thebibliography}{100}

\bibitem{Ykwon}
Y.~Kwon,
``Defy gravity in quantum field,''
Class.\ Quant. Grav.\ {\bf 19}, 9999 (2024) 
\end{thebibliography}

\endgroup
%%%%%%%%%%%%%%%%%%%%%%%%%%%%%%%%
%---(Alternatively)----%%%%%%%%%%%%%%%
%\bibliographystyle{unsrt}
%\bibliography{biblio}
%%%%%%%%%%%%%%%%%%%%%%%%%%%%%%







\end{document}




